% assets/w04-sensor-layout.svg — WAM-V 센서 배치 (실측 좌표 + 시야)
%
%   bash _tools/tikz2svg.sh assets/src/w04-sensor-layout.tex assets/w04-sensor-layout.svg
%
% 요점 두 가지
%   1) 센서는 전부 앞쪽에 몰려 있다 -> 후방이 사각지대다
%   2) 좌표는 VRX 실행 중 TF 에서 잰 값이다 (추정 아님)
%
% 선체 좌표계 : x 전방, y 좌현, z 상방 (ENU 계열)
% 화면 : x 오른쪽, y 위쪽.  1 m = 15 mm

\documentclass[border=6pt]{standalone}
% capstone-style.tex — 이 볼트의 TikZ 그림이 공유하는 스타일
%
% 쓰는 법:  % capstone-style.tex — 이 볼트의 TikZ 그림이 공유하는 스타일
%
% 쓰는 법:  % capstone-style.tex — 이 볼트의 TikZ 그림이 공유하는 스타일
%
% 쓰는 법:  \input{capstone-style}   (assets/src/ 안에서)
% 빌드:     bash _tools/tikz2svg.sh assets/src/w02-node-topic.tex assets/w02-node-topic.svg
%
% 왜 TikZ 인가
%   손으로 SVG 좌표를 쓰면 화살촉·선 굵기·글자 크기가 그림마다 미묘하게 달라진다.
%   그것이 "자료가 어설퍼 보인다"의 정체다. 스타일을 한 번만 정하고 상대좌표로
%   배치하면 좌표를 고쳐도 그림이 무너지지 않는다.
%   수식은 본문과 같은 TeX 글꼴로 나온다.
%
% 한글
%   ko.TeX(DVI 경로)를 쓴다. XeLaTeX 은 TikZ 그래픽을 PDF special 로 내보내는데
%   dvisvgm 이 그것을 받으려면 Ghostscript < 10.01 이 필요하다. 이 PC 에는
%   10.07 이 깔려 있어 그 경로가 막힌다. latex -> DVI -> dvisvgm 이 유일하게
%   동작하는 조합이다.

\usepackage[hangul]{kotex}
\usepackage{tikz}
\usepackage{amsmath}
\usepackage{xcolor}
\usetikzlibrary{arrows.meta, positioning, calc, fit, backgrounds, decorations.pathmorphing}

% ---- 색 --------------------------------------------------------------------
% 강조는 하나만. 나머지는 검정. 흑백 인쇄에서도 읽혀야 한다.
\definecolor{ink}{HTML}{1A1A1A}
\definecolor{soft}{HTML}{666666}
\definecolor{accent}{HTML}{7C3AED}
\definecolor{warn}{HTML}{B91C1C}

% 계층 색 — Simulink 모델의 블록 색과 같은 규칙을 쓴다
\definecolor{guid}{HTML}{1D4ED8}    % 유도
\definecolor{ctrl}{HTML}{EA580C}    % 제어
\definecolor{thr}{HTML}{D97706}     % 추진기
\definecolor{plant}{HTML}{16A34A}   % 운동모델
\definecolor{nav}{HTML}{7C3AED}     % 항법 · 신호처리
\definecolor{pathc}{HTML}{0D9488}   % 계획 경로 · 항적
\definecolor{errc}{HTML}{C2410C}    % 오차
\definecolor{flow}{HTML}{0369A1}    % 조류 · 바람

% ---- 선과 화살촉 -----------------------------------------------------------
\tikzset{
  sig/.style   = {ink, line width=0.5pt, -{Stealth[length=4pt, width=3pt]}},
  wire/.style  = {ink, line width=0.5pt},
  fb/.style    = {ink, line width=0.5pt, densely dashed,
                  -{Stealth[length=4pt, width=3pt]}},
  asig/.style  = {accent, line width=0.5pt, densely dashed,
                  -{Stealth[length=4pt, width=3pt]}},
  % 블록 — 흰 바탕, 직각, 검은 테두리. 색은 테두리에만 준다
  blk/.style   = {draw=ink, line width=0.55pt, fill=white,
                  minimum width=18mm, minimum height=8.5mm, inner sep=2pt, font=\small},
  gblk/.style  = {blk, draw=guid},
  cblk/.style  = {blk, draw=ctrl},
  tblk/.style  = {blk, draw=thr},
  pblk/.style  = {blk, draw=plant},
  nblk/.style  = {blk, draw=nav},
  % 합산점 · 분기점
  sum/.style   = {draw=ink, line width=0.55pt, fill=white, circle,
                  inner sep=0pt, minimum size=4.6mm, font=\scriptsize},
  dot/.style   = {ink, fill=ink, circle, inner sep=0pt, minimum size=1.5mm},
  % 글자
  sname/.style = {font=\small, inner sep=1.5pt},
  note/.style  = {font=\scriptsize, soft, inner sep=1.5pt},
  wnote/.style = {font=\scriptsize, warn, inner sep=1.5pt},
  anote/.style = {font=\scriptsize, accent, inner sep=1.5pt},
  axis/.style  = {ink, line width=0.5pt},
  tick/.style  = {font=\scriptsize, soft},
  % 묶음 상자 — 서브시스템 표시
  grp/.style   = {draw=soft, line width=0.4pt, densely dotted, rounded corners=1mm,
                  inner sep=3mm},
}
   (assets/src/ 안에서)
% 빌드:     bash _tools/tikz2svg.sh assets/src/w02-node-topic.tex assets/w02-node-topic.svg
%
% 왜 TikZ 인가
%   손으로 SVG 좌표를 쓰면 화살촉·선 굵기·글자 크기가 그림마다 미묘하게 달라진다.
%   그것이 "자료가 어설퍼 보인다"의 정체다. 스타일을 한 번만 정하고 상대좌표로
%   배치하면 좌표를 고쳐도 그림이 무너지지 않는다.
%   수식은 본문과 같은 TeX 글꼴로 나온다.
%
% 한글
%   ko.TeX(DVI 경로)를 쓴다. XeLaTeX 은 TikZ 그래픽을 PDF special 로 내보내는데
%   dvisvgm 이 그것을 받으려면 Ghostscript < 10.01 이 필요하다. 이 PC 에는
%   10.07 이 깔려 있어 그 경로가 막힌다. latex -> DVI -> dvisvgm 이 유일하게
%   동작하는 조합이다.

\usepackage[hangul]{kotex}
\usepackage{tikz}
\usepackage{amsmath}
\usepackage{xcolor}
\usetikzlibrary{arrows.meta, positioning, calc, fit, backgrounds, decorations.pathmorphing}

% ---- 색 --------------------------------------------------------------------
% 강조는 하나만. 나머지는 검정. 흑백 인쇄에서도 읽혀야 한다.
\definecolor{ink}{HTML}{1A1A1A}
\definecolor{soft}{HTML}{666666}
\definecolor{accent}{HTML}{7C3AED}
\definecolor{warn}{HTML}{B91C1C}

% 계층 색 — Simulink 모델의 블록 색과 같은 규칙을 쓴다
\definecolor{guid}{HTML}{1D4ED8}    % 유도
\definecolor{ctrl}{HTML}{EA580C}    % 제어
\definecolor{thr}{HTML}{D97706}     % 추진기
\definecolor{plant}{HTML}{16A34A}   % 운동모델
\definecolor{nav}{HTML}{7C3AED}     % 항법 · 신호처리
\definecolor{pathc}{HTML}{0D9488}   % 계획 경로 · 항적
\definecolor{errc}{HTML}{C2410C}    % 오차
\definecolor{flow}{HTML}{0369A1}    % 조류 · 바람

% ---- 선과 화살촉 -----------------------------------------------------------
\tikzset{
  sig/.style   = {ink, line width=0.5pt, -{Stealth[length=4pt, width=3pt]}},
  wire/.style  = {ink, line width=0.5pt},
  fb/.style    = {ink, line width=0.5pt, densely dashed,
                  -{Stealth[length=4pt, width=3pt]}},
  asig/.style  = {accent, line width=0.5pt, densely dashed,
                  -{Stealth[length=4pt, width=3pt]}},
  % 블록 — 흰 바탕, 직각, 검은 테두리. 색은 테두리에만 준다
  blk/.style   = {draw=ink, line width=0.55pt, fill=white,
                  minimum width=18mm, minimum height=8.5mm, inner sep=2pt, font=\small},
  gblk/.style  = {blk, draw=guid},
  cblk/.style  = {blk, draw=ctrl},
  tblk/.style  = {blk, draw=thr},
  pblk/.style  = {blk, draw=plant},
  nblk/.style  = {blk, draw=nav},
  % 합산점 · 분기점
  sum/.style   = {draw=ink, line width=0.55pt, fill=white, circle,
                  inner sep=0pt, minimum size=4.6mm, font=\scriptsize},
  dot/.style   = {ink, fill=ink, circle, inner sep=0pt, minimum size=1.5mm},
  % 글자
  sname/.style = {font=\small, inner sep=1.5pt},
  note/.style  = {font=\scriptsize, soft, inner sep=1.5pt},
  wnote/.style = {font=\scriptsize, warn, inner sep=1.5pt},
  anote/.style = {font=\scriptsize, accent, inner sep=1.5pt},
  axis/.style  = {ink, line width=0.5pt},
  tick/.style  = {font=\scriptsize, soft},
  % 묶음 상자 — 서브시스템 표시
  grp/.style   = {draw=soft, line width=0.4pt, densely dotted, rounded corners=1mm,
                  inner sep=3mm},
}
   (assets/src/ 안에서)
% 빌드:     bash _tools/tikz2svg.sh assets/src/w02-node-topic.tex assets/w02-node-topic.svg
%
% 왜 TikZ 인가
%   손으로 SVG 좌표를 쓰면 화살촉·선 굵기·글자 크기가 그림마다 미묘하게 달라진다.
%   그것이 "자료가 어설퍼 보인다"의 정체다. 스타일을 한 번만 정하고 상대좌표로
%   배치하면 좌표를 고쳐도 그림이 무너지지 않는다.
%   수식은 본문과 같은 TeX 글꼴로 나온다.
%
% 한글
%   ko.TeX(DVI 경로)를 쓴다. XeLaTeX 은 TikZ 그래픽을 PDF special 로 내보내는데
%   dvisvgm 이 그것을 받으려면 Ghostscript < 10.01 이 필요하다. 이 PC 에는
%   10.07 이 깔려 있어 그 경로가 막힌다. latex -> DVI -> dvisvgm 이 유일하게
%   동작하는 조합이다.

\usepackage[hangul]{kotex}
\usepackage{tikz}
\usepackage{amsmath}
\usepackage{xcolor}
\usetikzlibrary{arrows.meta, positioning, calc, fit, backgrounds, decorations.pathmorphing}

% ---- 색 --------------------------------------------------------------------
% 강조는 하나만. 나머지는 검정. 흑백 인쇄에서도 읽혀야 한다.
\definecolor{ink}{HTML}{1A1A1A}
\definecolor{soft}{HTML}{666666}
\definecolor{accent}{HTML}{7C3AED}
\definecolor{warn}{HTML}{B91C1C}

% 계층 색 — Simulink 모델의 블록 색과 같은 규칙을 쓴다
\definecolor{guid}{HTML}{1D4ED8}    % 유도
\definecolor{ctrl}{HTML}{EA580C}    % 제어
\definecolor{thr}{HTML}{D97706}     % 추진기
\definecolor{plant}{HTML}{16A34A}   % 운동모델
\definecolor{nav}{HTML}{7C3AED}     % 항법 · 신호처리
\definecolor{pathc}{HTML}{0D9488}   % 계획 경로 · 항적
\definecolor{errc}{HTML}{C2410C}    % 오차
\definecolor{flow}{HTML}{0369A1}    % 조류 · 바람

% ---- 선과 화살촉 -----------------------------------------------------------
\tikzset{
  sig/.style   = {ink, line width=0.5pt, -{Stealth[length=4pt, width=3pt]}},
  wire/.style  = {ink, line width=0.5pt},
  fb/.style    = {ink, line width=0.5pt, densely dashed,
                  -{Stealth[length=4pt, width=3pt]}},
  asig/.style  = {accent, line width=0.5pt, densely dashed,
                  -{Stealth[length=4pt, width=3pt]}},
  % 블록 — 흰 바탕, 직각, 검은 테두리. 색은 테두리에만 준다
  blk/.style   = {draw=ink, line width=0.55pt, fill=white,
                  minimum width=18mm, minimum height=8.5mm, inner sep=2pt, font=\small},
  gblk/.style  = {blk, draw=guid},
  cblk/.style  = {blk, draw=ctrl},
  tblk/.style  = {blk, draw=thr},
  pblk/.style  = {blk, draw=plant},
  nblk/.style  = {blk, draw=nav},
  % 합산점 · 분기점
  sum/.style   = {draw=ink, line width=0.55pt, fill=white, circle,
                  inner sep=0pt, minimum size=4.6mm, font=\scriptsize},
  dot/.style   = {ink, fill=ink, circle, inner sep=0pt, minimum size=1.5mm},
  % 글자
  sname/.style = {font=\small, inner sep=1.5pt},
  note/.style  = {font=\scriptsize, soft, inner sep=1.5pt},
  wnote/.style = {font=\scriptsize, warn, inner sep=1.5pt},
  anote/.style = {font=\scriptsize, accent, inner sep=1.5pt},
  axis/.style  = {ink, line width=0.5pt},
  tick/.style  = {font=\scriptsize, soft},
  % 묶음 상자 — 서브시스템 표시
  grp/.style   = {draw=soft, line width=0.4pt, densely dotted, rounded corners=1mm,
                  inner sep=3mm},
}


\begin{document}
\begin{tikzpicture}[x=1mm, y=1mm]

\def\S{15}   % 1 m

% ===================== 평면도 =====================
\begin{scope}
  \node[font=\small\bfseries, anchor=west] at (-46,26) {평면도 — 위에서 본 모습};

  % 시야 : 전방 부채꼴
  \fill[pathc, opacity=0.10] (0,0) -- (45:34mm) arc (45:-45:34mm) -- cycle;
  \node[note, pathc, anchor=west] at (26,-13) {카메라 · LiDAR 시야};

  % 사각지대 : 후방 부채꼴
  \fill[warn, opacity=0.10] (0,0) -- (150:26mm) arc (150:210:26mm) -- cycle;
  \node[note, warn, align=center] at (-30,0) {후방\\ 사각지대};

  % 선체 두 개 (y = +-1.027 m)
  \draw[flow, line width=0.7pt, rounded corners=2mm]
       (-2.374*\S-9,  1.027*\S-3) rectangle (0.9*\S+9,  1.027*\S+3);
  \draw[flow, line width=0.7pt, rounded corners=2mm]
       (-2.374*\S-9, -1.027*\S-3) rectangle (0.9*\S+9, -1.027*\S+3);
  \node[note, flow, anchor=west] at (0.9*\S+11,  1.027*\S) {좌현 $y=+1.03$};
  \node[note, flow, anchor=west] at (0.9*\S+11, -1.027*\S) {우현 $y=-1.03$};

  % 기준점과 축
  \fill[ink] (0,0) circle (1.1);
  \node[note, anchor=north east] at (-1,-1) {\texttt{base\_link}};
  \draw[axis, -{Stealth[length=4pt,width=3pt]}] (0,0) -- (16,0);
  \node[note, anchor=north] at (17,-1) {$x$ 선수};
  \draw[axis, -{Stealth[length=4pt,width=3pt]}] (0,0) -- (0,12);
  \node[note, anchor=east] at (-0.5,12) {$y$};

  % 센서
  \fill[errc] (0.700*\S,0) circle (1.6);
  \node[note, errc, anchor=south] at (0.700*\S-2,5) {3D LiDAR};
  \fill[thr]  (0.700*\S,-3.4) circle (1.3);
  \node[note, thr, anchor=north west] at (0.700*\S+2,-4) {Livox};
  \fill[pathc] (0.750*\S, 0.100*\S) circle (1.2);
  \fill[pathc] (0.750*\S,-0.100*\S) circle (1.2);
  \node[note, pathc, anchor=west] at (0.750*\S+4,0.100*\S+4) {전방 카메라 2};
  \fill[pathc] (0.500*\S,-0.450*\S) circle (1.2);
  \node[note, pathc, anchor=north] at (0.500*\S,-0.450*\S-2) {우현 카메라};
  \fill[flow] (-0.850*\S,0) circle (1.5);
  \node[note, flow, anchor=south] at (-0.850*\S,2) {GPS};
  \fill[nav]  (0.300*\S,-0.200*\S) circle (1.3);
  \node[note, nav, anchor=north east] at (0.300*\S-1,-0.200*\S-1) {IMU};

  % 추진기
  \draw[thr, line width=0.7pt] (-2.374*\S-4,  1.027*\S-3) rectangle (-2.374*\S+4,  1.027*\S+3);
  \draw[thr, line width=0.7pt] (-2.374*\S-4, -1.027*\S-3) rectangle (-2.374*\S+4, -1.027*\S+3);
  \node[note, thr, anchor=east] at (-2.374*\S-6, 1.027*\S) {추진기 $x=-2.37$};
\end{scope}

% ===================== 측면도 =====================
\begin{scope}[shift={(0,-70)}]
  \node[font=\small\bfseries, anchor=west] at (-46,30) {측면도 — 옆에서 본 모습};

  \draw[flow, densely dashed] (-46,0) -- (46,0);
  \node[note, flow, anchor=west] at (47,0) {수면};
  \draw[flow, line width=0.7pt, rounded corners=2mm]
       (-2.374*\S-9,-3) rectangle (0.9*\S+9,3);

  \fill[ink] (0,0) circle (1.1);
  \node[note, anchor=north] at (0,-2) {\texttt{base\_link}};

  \draw[soft, line width=0.5pt] (0.700*\S,0) -- (0.700*\S,1.800*\S);
  \fill[errc] (0.700*\S,1.800*\S) circle (1.6);
  \node[note, errc, anchor=west] at (0.700*\S+3,1.800*\S) {3D LiDAR  $z=1.80$};
  \fill[pathc] (0.750*\S,1.500*\S) circle (1.2);
  \node[note, pathc, anchor=west] at (0.750*\S+3,1.500*\S) {카메라  $z=1.50$};
  \fill[thr] (0.700*\S,0.710*\S) circle (1.3);
  \node[note, thr, anchor=west] at (0.700*\S+3,0.710*\S) {Livox  $z=0.71$};

  \draw[soft, line width=0.5pt] (-0.850*\S,0) -- (-0.850*\S,1.300*\S);
  \fill[flow] (-0.850*\S,1.300*\S) circle (1.5);
  \node[note, flow, anchor=east] at (-0.850*\S-3,1.300*\S) {GPS  $z=1.30$};

  \fill[thr] (-2.652*\S,-0.191*\S) circle (1.3);
  \node[note, thr, anchor=east] at (-2.652*\S-3,-0.191*\S) {프로펠러 $z=-0.19$};
\end{scope}

% ===================== 주의 =====================
\node[wnote, anchor=west, text width=104mm] at (-46,-100)
     {센서가 전부 앞쪽에 몰려 있어 \textbf{후방이 사각지대}다.
      좌표는 VRX 실행 중 TF 에서 측정한 값이며(2026-09-06), 4주차에 직접 바꿔 본다.
      이 그림의 $y$ 는 \textbf{좌현이 양(+)} 이다 — 7주차 이후 NED 수식과 부호가 반대다};

\end{tikzpicture}
\end{document}
